\documentclass[11pt,letterpaper]{article}
\usepackage[top=0.7in,bottom=0.7in,left=0.75in,right=0.75in]{geometry}
\usepackage[T1]{fontenc}
\usepackage[utf8]{inputenc}
\usepackage{lmodern}
\usepackage{microtype}
\usepackage{enumitem}
\usepackage[hidelinks]{hyperref}
\usepackage{titlesec}
\setlength{\parindent}{0pt}
\setlength{\parskip}{0pt}
\pagestyle{empty}
\titleformat{\section}{\normalfont\bfseries}{}{0em}{\uppercase}[\vspace{-4pt}\hrule\vspace{4pt}]
\titlespacing{\section}{0pt}{8pt}{5pt}
\setlist[itemize]{leftmargin=1.2em,label=--,itemsep=0pt,topsep=2pt,parsep=0pt}
\begin{document}
\begin{center}
{\fontsize{16}{19}\selectfont\textbf{\MakeUppercase{Diego Castillo Pineda}}}\\[4pt]
{\small La Florida, Santiago, Chile \quad|\quad +56 9 94162547 \quad|\quad \href{mailto:diego.castillop11@gmail.com}{diego.castillop11@gmail.com}\\
\href{https://https://www.linkedin.com/in/diego-castillo-pineda/}{https://www.linkedin.com/in/diego-castillo-pineda/} \quad|\quad \href{https://https://github.com/diegocastillop11-cloud}{https://github.com/diegocastillop11-cloud}}
\end{center}

\section{Resumen Profesional}
Analista de Datos BI con experiencia práctica en SQL Server, Azure SQL, Oracle y Power BI, enfocado en la optimización de bases de datos transaccionales de alta concurrencia y en la automatización de procesos ETL. He diseñado stored procedures que procesaron decenas de miles de transacciones por hora en entornos productivos críticos y desarrollado consultas dinámicas PIVOT para reportería avanzada. Mi trayectoria incluye administración de bases de datos, análisis funcional, integración de sistemas mediante JSON/XML, y conocimientos en Python, C\# y herramientas modernas de visualización de datos como Power BI. Busco aplicar mi experiencia técnica y analítica en Step Therapy para contribuir al desarrollo de soluciones de Business Intelligence escalables y basadas en datos.

\section{Experiencia Profesional}

\textbf{Punto Ticket} \hfill Santiago, Chile\\
\textit{Analista de Base de Datos Senior} \hfill Nov. 2019 – Feb. 2026
\begin{itemize}
  \item Diseñé y optimicé stored procedures en SQL Server para procesar decenas de miles de transacciones concurrentes por hora durante eventos masivos, manteniendo tiempos de respuesta dentro de umbrales operativos críticos y garantizando disponibilidad del 99,9\% en períodos de máxima demanda
  \item Desarrollé consultas dinámicas con PIVOT y funciones avanzadas de T-SQL para reportería ejecutiva, reemplazando procesos manuales y reduciendo en un 70\% el tiempo de generación de informes mensuales
  \item Automaticé procesos masivos mediante scripts T-SQL y configuraciones de SQL Agent Jobs, disminuyendo errores operativos en un 80\% y liberando más de 15 horas mensuales de trabajo manual del equipo
  \item Implementé integraciones de datos entre sistemas utilizando JSON, XML y APIs REST, facilitando la sincronización en tiempo real de información transaccional entre plataformas de venta y sistemas de contabilidad
  \item Participé activamente en la migración de bases de datos on-premise hacia Microsoft Azure (Azure SQL Database), colaborando en la definición de arquitectura cloud, pruebas de carga y ajuste de rendimiento post-migración
  \item Administré bases de datos SQL Server en entornos de alta concurrencia, aplicando técnicas de indexación, particionamiento y monitoreo proactivo con SQL Server Management Studio y Azure Monitor para identificar y resolver cuellos de botella antes de impactos productivos
\end{itemize}
\vspace{2pt}
\textbf{Imperial S.A.} \hfill Santiago, Chile\\
\textit{Analista Funcional} \hfill Jun. 2017 – Nov. 2019
\begin{itemize}
  \item Realicé levantamiento y análisis de requerimientos funcionales junto a usuarios clave de distintas áreas operativas, traduciendo necesidades de negocio en especificaciones técnicas documentadas para equipos de desarrollo y soporte
  \item Brindé soporte funcional y técnico a sistemas ERP (Oracle, SAP) y aplicaciones en Visual Basic 6, resolviendo incidencias de datos mediante consultas SQL y scripts de corrección en ambientes productivos y de prueba
  \item Configuré y parametricé módulos de sistemas según especificaciones funcionales, elaborando documentación completa (instructivos, flujos de proceso, manuales de usuario) que redujo en un 50\% las consultas recurrentes de soporte
  \item Capacité a usuarios finales en el uso de nuevas funcionalidades y cambios de sistemas, facilitando la adopción tecnológica y mejorando la eficiencia operativa en procesos de facturación y gestión de inventarios
\end{itemize}
\vspace{2pt}
\textbf{OB GROUP PARK} \hfill Santiago, Chile\\
\textit{Analista de Sistemas y TI} \hfill Mayo 2014 – Abril 2017
\begin{itemize}
  \item Lideré proyectos de mejora de procesos contables y de gestión en el sector retail, coordinando la integración de nuevas tecnologías (sistemas POS, herramientas de reportería) y asegurando compatibilidad entre plataformas legacy y modernas
  \item Colaboré con áreas de contabilidad, finanzas y operaciones para garantizar la continuidad operativa durante implementaciones tecnológicas, documentando flujos de datos y realizando validaciones de integridad entre sistemas
  \item Implementé soluciones técnicas para automatizar reportes de ventas diarias, reduciendo el tiempo de cierre contable mensual en un 30\% mediante la integración de datos desde múltiples fuentes (Excel, SQL, archivos CSV)
\end{itemize}
\vspace{2pt}
\textbf{Hewlett-Packard Company} \hfill Santiago, Chile\\
\textit{Agente de Mesa de Ayuda Nivel 1} \hfill Ago. 2013 – Mayo 2014
\begin{itemize}
  \item Proporcioné soporte técnico de primer nivel a usuarios finales, diagnosticando y resolviendo incidencias de hardware, software y conectividad, desarrollando habilidades de comunicación técnica y gestión de incidentes que guían mi trabajo actual
\end{itemize}

\section{Proyectos Destacados}

\textbf{Plataforma de Análisis Predictivo con IA y React} \hfill 2024
\begin{itemize}
  \item Desarrollé una aplicación full-stack para análisis de datos utilizando React (TypeScript) en el frontend, Node.js y Python (FastAPI) en el backend, conectada a PostgreSQL y MongoDB para almacenamiento estructurado y no estructurado
  \item Integré modelos de IA mediante prompt engineering con Claude y Gemini API para generación automática de insights y recomendaciones a partir de datasets de ventas, logrando respuestas en menos de 3 segundos para conjuntos de 50.000+ registros
  \item Implementé autenticación segura con JWT, encriptación de datos sensibles en tránsito (TLS) y en reposo (bcrypt), y diseñé una UI/UX intuitiva con componentes reutilizables y dashboards interactivos que mejoraron la experiencia de usuario en un 40\% según feedback de pruebas
\end{itemize}

\section{Habilidades T\'{e}cnicas}
\textbf{Datos \& Cloud:} SQL Server (Experto), Azure SQL, Oracle, PostgreSQL, MongoDB, Microsoft Azure \\
\textbf{Desarrollo \& IA:} Python, C\#, Node.js, Go, React, Prompt Engineering (Claude / Gemini) \\
\textbf{Frontend \& Diseño:} TypeScript, JavaScript (ES6+), HTML5/CSS3, UI/UX Design \\
\textbf{Herramientas \& Metodologías:} Git, Docker, Power BI, Excel Avanzado, Scrum/Kanban

\section{Formaci\'{o}n Acad\'{e}mica}
\textbf{Business Analytics \& Data Science}, Universidad de Chile \hfill 2024 \\
\textbf{Especialización en Optimización de Procesos y Administración SQL Server}, Universidad de Chile \hfill 2023 \\
\textbf{Desarrollo de Apps Móviles}, Universidad Complutense de Madrid \hfill 2017 \\
\textbf{Analista Programador}, INACAP \hfill 2013

\end{document}